\chapter{Exponentials and the Number e}

In Chapter~2, we explored algebraic functions where the variable $x$ serves as a
base raised to fixed powers, like $x^2$ or $\sqrt{x}$. But what happens when that
relationship is turned inside out -- when a fixed base is raised to a variable
exponent $x$, producing curves like $2^x$ that multiply rather than add?

Among all possible bases, one number stands out across calculus: $e \approx 2.71828$. People often meet $e$ as a button on a calculator and are told it is the ``natural'' base without seeing why a strange decimal near $2.718$ could be more natural than $2$ or $10$. In this chapter we will see where $e$ comes from -- and why compound interest, tangent slopes, and infinite series all lead straight to the same number.

We start with the general family of exponential functions.

\section{What an Exponential Really Is}

\subsection{From repeated multiplication to a smooth curve}

An \textbf{exponential function} is a function of the form $f(x)=a^x$, where the base $a$ is a fixed positive number and the input variable $x$ sits in the exponent.

You first met powers as repeated multiplication. If $a$ is a positive number,
\[
a^1 = a, \qquad a^2 = a\cdot a, \qquad a^3 = a\cdot a\cdot a, \qquad a^n = \underbrace{a\cdot a\cdots a}_{n\text{ factors}}.
\]
This is airtight when $n$ is a positive whole number. You cannot multiply $a$ by itself ``$\tfrac{1}{2}$ times'' or ``$-3$ times'' or ``$\pi$ times'' by counting factors. So if we want a function $f(x)=a^x$ defined for \emph{every} real number $x$ -- a smooth curve, rather than merely a discrete scatter of points over the integers -- we have to \emph{extend} the idea. The remarkable fact is that there is exactly one consistent way to do it.

The whole extension is forced by a single property we refuse to give up:
\begin{equation}
a^{x+y} = a^x \, a^y. \label{ch03:law}
\end{equation}
This is the law of exponents, and over the integers it is just bookkeeping: $a^{2+3}=a^2 a^3$ because $5$ factors is $2$ factors followed by $3$ factors. We will \emph{insist} that this law keeps holding for all real exponents, and that single insistence pins down everything else.

\begin{intuitionbox}
Think of $a^x$ as a machine that turns \emph{addition into multiplication}. Slide the input $x$ along by a fixed step $\Delta x$, and the output does not change by a fixed \emph{amount} -- it changes by a fixed \emph{factor}. Walk one unit to the right and you multiply by $a$. Walk another unit and you multiply by $a$ again. Exponentials are the functions for which equal strides cause equal \emph{percentage} changes. That is the whole personality of the thing.
\end{intuitionbox}

Let us see how \eqref{ch03:law} forces the rest of the definition, step by step.

\textbf{Zero.} Put $y=0$ in \eqref{ch03:law}: $a^{x+0}=a^x a^0$, i.e. $a^x = a^x\, a^0$. Dividing by $a^x$ (which is positive, never zero) gives
\[
a^0 = 1.
\]
We did not \emph{decree} that $a^0=1$ to be tidy; the exponent law leaves us no choice.

\textbf{Negative exponents.} Put $y=-x$: $a^{x+(-x)} = a^0 = 1$, so $a^x a^{-x}=1$, hence
\[
a^{-x} = \frac{1}{a^x}.
\]
Again, forced. A negative exponent must mean a reciprocal, because that is the only thing consistent with the law.

\textbf{Roots and fractions.} What could $a^{1/2}$ be? By the law, $a^{1/2}\cdot a^{1/2} = a^{1/2+1/2}=a^1=a$. So $a^{1/2}$ is a number whose square is $a$: it is $\sqrt{a}$. More generally $a^{1/n}=\sqrt[n]{a}$, and $a^{m/n}=\big(\sqrt[n]{a}\big)^m$. Every rational exponent is now determined.

\textbf{Irrational exponents.} This is the only delicate step, and it is the place where the integers quietly hand the baton to calculus. What is $2^{\pi}$? We cannot reach $\pi$ by counting or by fractions exactly. But we can \emph{sneak up on it}. The number $\pi = 3.14159\ldots$ is the limit of the rational numbers
\[
3,\quad 3.1,\quad 3.14,\quad 3.141,\quad 3.1415,\ \ldots
\]
and each of those has a perfectly good power: $2^3, 2^{3.1}, 2^{3.14}, \ldots$ These values march steadily toward a single limiting number, and \emph{that} limit is what we define $2^{\pi}$ to be.

\begin{happeningbox}
When you write $2^{\pi}$, nothing is being ``multiplied $\pi$ times.'' What is really happening is a \emph{limit}. We define
\[
a^x \;=\; \lim_{\substack{r\to x \\ r\ \text{rational}}} a^{r},
\]
the value the rational powers approach as their exponent closes in on $x$. The reason this works -- the reason the rational powers do not wander off but settle on one number -- is that $a^r$ is \emph{monotonic and continuous} in $r$, a fact guaranteed by the completeness of the real numbers (the property that guarantees the number line has no gaps or missing points). Here, \emph{monotonic} simply means the values head strictly in one direction---always increasing for $a>1$, never turning around or fluctuating. So the exponential is the first place in this book where the smoothness of the real line is doing real labor under the surface. Calculus is already here, hiding inside an innocent-looking $2^x$.
\end{happeningbox}

\begin{figure}[htbp]
\centering
\begin{tikzpicture}
\begin{axis}[
  width=11.5cm, height=6.8cm,
  title style={font=\small},
  xlabel={$x$}, ylabel={$y$},
  domain=-1.5:3.2, samples=120,
  ymin=-0.5, ymax=8.5, xmin=-1.6, xmax=3.4,
  axis lines=middle,
]
\addplot[curveblue, very thick] {2^x};
\node[curveblue, font=\small, right] at (axis cs:2.8, 7.5) {$y=2^x$};
\addplot[only marks, sourceblue, mark=*, mark size=2.5pt] coordinates {
  (-1, 0.5) (0, 1) (1, 2) (2, 4) (3, 8)
};
\addplot[only marks, quotientorange, mark=*, mark size=2.0pt] coordinates {
  (-0.5, 0.7071) (0.5, 1.4142) (1.5, 2.8284) (2.5, 5.6569)
};
\end{axis}
\end{tikzpicture}
\caption{From discrete multiplication to a continuous curve. Counting factors defines $2^x$ only at integer exponents (blue dots). Demanding that the exponent law $2^{x+y}=2^x 2^y$ hold everywhere forces the values at fractional exponents, like $2^{1/2}=\sqrt{2}$ and $2^{3/2}=\sqrt{8}$ (orange dots). Finally, real completeness fills every intermediate gap, weaving the points into a single unbroken curve.}
\label{ch03:fig-discrete-to-smooth}
\end{figure}

\subsection{The shape of the graph}

With $a^x$ now defined for all real $x$, let us look at it. Take $a>1$, say $a=2$. The function climbs, gently at first and then ferociously. Take $0<a<1$, say $a=\tfrac12$; then $a^x = (1/2)^x = 2^{-x}$ is just the mirror image, decaying instead of growing. And $a=1$ is the boring flat line $1^x=1$.

\begin{figure}[htbp]
\centering
\begin{tikzpicture}
\begin{axis}[
  width=11.5cm, height=6.8cm,
  title style={font=\small},
  xlabel={$x$}, ylabel={$y$},
  domain=-3:3, samples=120,
  ymin=-0.4, ymax=8.4, xmin=-3.4, xmax=3.4,
  axis lines=middle,
  xtick={-3, -2, -1, 1, 2, 3}, ytick={2, 4, 6, 8},
  legend style={at={(0.5,-0.18)}, anchor=north, legend columns=2},
  legend cell align=left,
]
\addplot[curveblue, very thick] {2^x};
\addlegendentry{$y=2^x$ (growth, $a>1$)}
\addplot[quotientorange, very thick] {(1/2)^x};
\addlegendentry{$y=(1/2)^x$ (decay, $0<a<1$)}
\addplot[structurepurple, thick, dashed] {3^x};
\addlegendentry{$y=3^x$ (steeper base)}
\addplot[gray, thick, dotted, domain=-3:3] {1};
\addlegendentry{$y=1^x=1$}
\node[circle,fill=black,inner sep=1.5pt] at (axis cs:0,1) {};
\node[font=\footnotesize, fill=white, inner sep=1pt, anchor=north west] at (axis cs:0.35, 0.68) {$(0,1)$};
\draw[->, >=stealth, black!70, thin] (axis cs:0.35, 0.68) -- (axis cs:0.06, 0.95);
\end{axis}
\end{tikzpicture}
\caption{Every exponential $y=a^x$ passes through $(0,1)$, stays strictly positive, and is either increasing ($a>1$) or decreasing ($0<a<1$). Growth and decay are mirror images: $(1/2)^x$ is $2^x$ reflected across the $y$-axis.}
\label{ch03:fig-exp-family}
\end{figure}

Three features of Figure~\ref{ch03:fig-exp-family} are worth saying out loud, because they define how every exponential behaves and we will lean on all three.

\begin{itemize}
\item \textbf{It is always positive.} $a^x>0$ for every $x$. The curve hugs the $x$-axis on one side but never touches or crosses it. (A positive number raised to any power stays positive; the $x$-axis is a \textbf{horizontal asymptote}---a flat line that the curve gets closer and closer to as $x \to -\infty$, without ever quite touching it.)
\item \textbf{It always passes through $(0,1)$.} Because $a^0=1$ for every base, the whole family pivots through the single point $(0,1)$.
\item \textbf{Equal steps multiply.} Step right by $1$ and the height multiplies by $a$. This is constant \emph{proportional} change, and it is what makes exponentials the law of anything that grows or shrinks at a rate set by its own current size: populations, money, radioactive atoms, cooling coffee, viral posts.
\end{itemize}

\begin{trapbox}
An exponential function is $a^x$: the variable is upstairs, in the exponent. A \emph{power} function is $x^a$: the variable is downstairs, in the base. They are completely different animals. $x^2$ is a gentle parabola; $2^x$ eventually leaves every power function in the dust. Confusing $2^x$ with $x^2$ -- or differentiating one as if it were the other -- is one of the most common early mistakes. Say it to yourself: ``exponential, variable on top.''
\end{trapbox}

\begin{walkbox}
Here is a famous gut-check. A lily pad doubles its area every day, and on day $30$ it covers the whole pond. On which day was the pond \emph{half} covered? Most people instinctively say day $15$. Walk with that for a minute and feel why the answer is day $29$ -- and why ``half'' happens at the very last step. That surprise -- where half of all the growth happens in the single final step -- is what makes exponential growth so deceptive.
\end{walkbox}

\begin{recoverybox}
\textbf{Pause and think: the curved midpoint.} Since $2^1 = 2$ and $2^2 = 4$, the straight arithmetic average of the two outputs is $\frac{2+4}{2} = 3$. But what is $2^{1.5} = 2^{3/2} = \sqrt{2^3} = \sqrt{8} \approx 2.828$? Why is $2^{1.5}$ strictly smaller than $3$?

Look at the curve $y=2^x$ in Figure~\ref{ch03:fig-discrete-to-smooth}. Because the curve bends upward (it is concave up), any straight chord connecting two points on the graph floats strictly \emph{above} the curve. The midpoint of the straight chord connecting $(1,2)$ and $(2,4)$ is at height $3$, while the actual point on the exponential curve sits beneath it at $\sqrt{8} \approx 2.828$.
\end{recoverybox}

\begin{applicationbox}
\textbf{Why nature loves exponentials.} Suppose a quantity changes at a rate proportional to how much of it there is right now. Money in an account earns interest in proportion to the balance. A bacterial colony makes new cells in proportion to the cells already present. A lump of uranium loses atoms in proportion to the atoms still intact. In symbols, the rate of change is $k$ times the amount, and -- as we will prove once we have derivatives in hand -- the only functions obeying that are exponentials, $A(t)=A_0\, a^{t}$ for a suitable base. The exponential is not one model among many. It is \emph{the} signature of self-proportional change, which is why it shows up everywhere from compound interest to carbon dating to the early phase of an epidemic.
\end{applicationbox}

\section{Logarithms: Running the Exponential Backwards}

\subsection{The question a logarithm answers}

Since $a^x$ (for $a>1$) is strictly increasing, it hits every positive height exactly once. That means it can be \emph{undone}. Given a positive output $y$, there is exactly one input $x$ with $a^x = y$. We give that input a name.

\begin{definition}[Logarithm]\label{ch03:def-log}
For a base $a>0$ with $a\neq 1$, the \textbf{logarithm} $\log_a y$ is the exponent you must raise $a$ to in order to get $y$:
\[
\log_a y = x \quad\Longleftrightarrow\quad a^x = y, \qquad (y>0).
\]
In words: $\log_a y$ answers the question ``$a$ to the \emph{what} gives $y$?''
\end{definition}

So $\log_2 8 = 3$ because $2^3=8$, and $\log_{10}1000=3$ because $10^3=1000$, and $\log_2 \tfrac{1}{4} = -2$ because $2^{-2}=\tfrac14$. Fundamentally, a logarithm is simply an exponent. The logarithm and the exponential are \textbf{inverse functions}: each undoes the other,
\[
a^{\log_a y} = y \quad (y>0), \qquad \log_a(a^x) = x \quad (\text{all } x).
\]
An inverse function is simply an ``undo'' operation: if the exponential takes an input $x$ and climbs to $y = a^x$, the logarithm takes that output $y$ and recovers the original $x$ that produced it. Geometrically, inverse functions are reflections across the line $y=x$ -- you swap the roles of input and output, which swaps the roles of the two axes.

\begin{recoverybox}
\textbf{Pause and think: comparing logarithms.} Without using a calculator, which is larger: $\log_2 10$ or $\log_3 10$?

Translate both expressions into the logarithm's core question: $\log_2 10$ asks ``$2$ raised to what power reaches $10$?'' while $\log_3 10$ asks ``$3$ raised to what power reaches $10$?'' Because the base $2$ is smaller than $3$, it needs more steps of doubling ($2^3 = 8 < 10 < 16 = 2^4$, giving about $3.32$) to reach $10$ than base $3$ needs of tripling ($3^2 = 9 < 10 < 27 = 3^3$, giving about $2.10$). Therefore, $\log_2 10 > \log_3 10$.
\end{recoverybox}

\begin{figure}[htbp]
\centering
\begin{tikzpicture}
\begin{axis}[
  width=10.5cm, height=8cm,
  title style={font=\small},
  xlabel={$x$}, ylabel={$y$},
  domain=-2.5:6, samples=120,
  ymin=-2.8, ymax=6.6, xmin=-2.8, xmax=6.6,
  axis lines=middle,
  xtick={-2, 2, 4}, ytick={-2, 2, 4, 6},
  legend style={at={(0.5,-0.15)}, anchor=north, legend columns=-1},
  legend cell align=left,
  axis equal image,
]
\addplot[curveblue, very thick, domain=-2.5:2.4] {2^x};
\addlegendentry{$y=2^x$}
\addplot[targetgreen, very thick, domain=0.06:6] {ln(x)/ln(2)};
\addlegendentry{$y=\log_2 x$}
\addplot[gray, thick, dashed, domain=-2.5:6] {x};
\addlegendentry{$y=x$}
\end{axis}
\end{tikzpicture}
\caption{The logarithm is the exponential reflected across the dashed line $y=x$. The point $(0,1)$ on the exponential becomes $(1,0)$ on the logarithm; the vertical asymptote of the log (at $x=0$) is the reflected image of the horizontal asymptote of the exponential.}
\label{ch03:fig-log-reflect}
\end{figure}

\subsection{Why logs turn multiplication into addition}

Every algebraic property of logarithms is just the exponent law \eqref{ch03:law} read backwards. This is worth seeing once, slowly, because it is the reason logarithms were invented and the reason they still matter.

\begin{whybox}
\textbf{The mechanism: a logarithm is an exponent, so adding exponents becomes adding logarithms.} Take two positive numbers $u$ and $v$. Write them in terms of the base: $u = a^p$ and $v = a^q$, which is to say $p=\log_a u$ and $q=\log_a v$. Now multiply:
\[
u\,v = a^p\, a^q = a^{p+q}
\]
by the exponent law. Read that last equation through the logarithm: the exponent that produces $uv$ is $p+q$. Therefore
\[
\boxed{\;\log_a(uv) = \log_a u + \log_a v.\;}
\]
\emph{Multiplication downstairs became addition upstairs.} The same one-line trick gives every other rule:
\[
\log_a\!\Big(\frac{u}{v}\Big) = \log_a u - \log_a v, \qquad \log_a(u^n) = n\,\log_a u.
\]
There is no separate ``log algebra'' to memorize. There is only the exponent law, seen from the other side.
\end{whybox}

\begin{applicationbox}
\textbf{How logarithms revolutionized astronomy, navigation, and spaceflight.}
This was not a mere theoretical curiosity -- it changed history. In 1614 John Napier published the first table of logarithms, and Henry Briggs soon recast them in base $10$. Before computers, multiplying two eight-digit numbers by hand was slow, exhausting, and error-prone. With a log table, an astronomer could \emph{look up} the logs, \emph{add} them (a simple sum), and look the answer back up. Multiplication had been demoted to addition. Kepler used logarithms to grind through the heavy astronomical arithmetic behind his laws of planetary motion; Laplace later remarked that logarithms, ``by shortening the labours, doubled the life of the astronomer.'' The slide rule -- the engineer's constant pocket companion until the 1970s -- consisted of two logarithmic rulers sliding past each other, \emph{adding physical lengths} to multiply numbers. Every Apollo mission to the Moon was flown with calculations checked on a device whose whole operating principle was $\log(uv)=\log u+\log v$.
\end{applicationbox}

\begin{focusbox}
\textbf{Changing base.} Any one logarithm can be rebuilt from any other by a constant factor. If $a^x=y$, take $\log_b$ of both sides: $x\log_b a = \log_b y$, so
\[
\log_a y = \frac{\log_b y}{\log_b a}.
\]
The shape of the log curve never really changes with the base -- only its vertical scale. That is exactly why, in a moment, we will be free to hunt for the \emph{one} base that makes calculus cleanest, knowing every other base is just a rescaling of it.
\end{focusbox}

\section{The Central Question: A Base That Grows At Its Own Rate}

Now we come to the heart of the chapter. We have a whole family of exponentials $a^x$, one for each base $a$. For calculus, the question that matters most is: \emph{how fast does $a^x$ change?} That is, what is the steepness of its graph?

In Chapter~1, we saw that the steepness of a smooth curve at any point can be inspected through a geometric magnifying glass: we pick a nearby point separated by a tiny horizontal step $h$, draw the straight line connecting them (called a \textbf{secant line}), and measure the average rate of change---the rise over the run.

We do not yet have the formal machinery of limits (which we will construct with full precision in Chapter~5) or the systematic rules of derivatives (which we will develop in Chapters 10 and 11). But exponential functions possess an algebraic symmetry so clean that we can preview their steepness right now, using only high-school algebra and secant slopes. When we reach Chapter 12, we will return to this exact calculation and establish it as a formal theorem of differential calculus. For now, let us see what the geometry reveals.

\subsection{Steepness and the starting slope of $a^x$}

Consider the secant slope between $(x, a^x)$ and a nearby point $(x+h, a^{x+h})$ on the curve $y=a^x$. The average rate of change over this tiny interval of width $h$ is the difference quotient:
\[
\frac{\Delta y}{\Delta x} \;=\; \frac{a^{x+h}-a^{x}}{h}.
\]
Using the law of exponents \eqref{ch03:law}, $a^{x+h} = a^x a^h$. Factoring out the common term $a^x$ gives:
\[
\frac{\Delta y}{\Delta x} \;=\; \frac{a^{x}a^{h}-a^{x}}{h}
       \;=\; a^{x}\cdot\frac{a^{h}-1}{h}.
\]
The average rate of change splits cleanly into two pieces: the height of the curve $a^x$, and a second factor $\frac{a^h-1}{h}$ that does not depend on the position $x$ at all!

As the step $h$ shrinks toward zero---an operation we denote with the standard limit notation $\lim_{h\to 0}$---the ratio $\frac{a^h-1}{h}$ settles toward a fixed number that depends only on the base $a$. Let us call this number $M(a)$:
\[
M(a) \;=\; \lim_{h\to 0}\frac{a^{h}-1}{h}.
\]
Because $a^x$ does not change as $h\to 0$, the instantaneous steepness of the curve at position $x$---traditionally written $\frac{\dd}{\dd x} a^x$ in the notation of Leibniz introduced in Chapter~1---is given by:
\begin{equation}
\frac{\dd}{\dd x}\,a^{x} \;=\; M(a)\cdot a^{x}. \label{ch03:propself}
\end{equation}

\begin{intuitionbox}
Read \eqref{ch03:propself} out loud: \emph{the slope of an exponential curve is a constant multiple of its current height.} Whatever the base, the rate of growth of $a^x$ is proportional to how high it has already climbed -- exactly the ``self-proportional'' behavior we said exponentials embody. The constant of proportionality $M(a)$ is the only thing that changes from base to base. And $M(a)$ has a beautifully simple meaning: setting $x=0$ in \eqref{ch03:propself} gives the slope at $x=0$ as $M(a)\cdot a^0 = M(a)$. So $M(a)$ is simply \emph{the slope of $a^x$ as it crosses the $y$-axis at the point $(0,1)$.}
\end{intuitionbox}

So every exponential is its own rate of change \emph{up to a constant factor} $M(a)$, the starting slope. And $M(a)$ is itself a limit of secant slopes through $(0,1)$: the quotient $(a^h-1)/h$ is the slope of the line joining $(0,1)$ to the nearby point $(h,a^h)$ on the curve. As $h$ shrinks, those secants pivot down onto the tangent at $(0,1)$, and their slopes settle on $M(a)$. Figure~\ref{ch03:fig-secant} shows the pivot in action.

\begin{figure}[htbp]
\centering
\begin{tikzpicture}
\begin{axis}[
  width=12cm, height=6.6cm,
  title style={font=\small},
  xlabel={$x$}, ylabel={$y$},
  domain=-0.4:1.3, samples=120,
  ymin=0.4, ymax=3.0, xmin=-0.45, xmax=1.45,
  axis lines=middle,
  xtick={-0.4, -0.2, 0.2, 0.4, 0.6, 0.8, 1.0, 1.2},
  ytick={0.5, 1.5, 2.0, 2.5, 3.0},
  legend pos=south east, legend cell align=left, legend style={font=\footnotesize},
]
\addplot[curveblue, very thick] {2^x};
\addlegendentry{$y=2^x$}
\addplot[gray, thick, dotted, domain=-0.4:1.3] {1 + 1.0*x};
\addlegendentry{secant, $h=1$ (slope $1.00$)}
\addplot[quotientorange, thick, dashed, domain=-0.4:1.3] {1 + 0.82843*x};
\addlegendentry{secant, $h=0.5$ (slope $0.83$)}
\addplot[tangentred, very thick, domain=-0.4:1.3] {1 + 0.6931*x};
\addlegendentry{tangent, slope $M(2)\approx0.69$}
\node[circle,fill=black,inner sep=1.2pt] at (axis cs:0,1){};
\node[anchor=south east, font=\footnotesize] at (axis cs:-0.03,1.05) {$(0,1)$};
\node[circle,fill=quotientorange,inner sep=1.2pt] at (axis cs:0.5,1.41421){};
\node[circle,fill=gray,inner sep=1.2pt] at (axis cs:1,2){};
\end{axis}
\end{tikzpicture}
\caption{The starting slope $M(a)$ as a limit of secants, shown for $a=2$. Each secant joins $(0,1)$ to a point $(h,2^h)$ on the curve; as $h\to 0$ the secants pivot down onto the tangent line at $(0,1)$, whose slope is $M(2)\approx 0.69$. The number $M(a)$ is the slope this pivoting line lands on.}
\label{ch03:fig-secant}
\end{figure}

This raises an irresistible question.

\begin{focusbox}
\textbf{The central question of this chapter.}
\[
\text{Is there a base } a \text{ for which the fudge factor is exactly } 1?
\]
That is, a base whose curve $a^x$ crosses $(0,1)$ at slope exactly $1$, so that $M(a)=1$ and the derivative of $a^x$ is $a^x$ \emph{on the nose}, with no leftover constant? Such a function would be its own rate of change -- the purest possible expression of self-proportional growth. \emph{That base, if it exists, is $e$.}
\end{focusbox}

\subsection{Cornering the base by measuring starting slopes}

Let us actually measure $M(a)$ for a few bases and watch the special base reveal itself. Numerically, with a small $h$, $M(a)\approx (a^h-1)/h$. Taking $h=0.0001$:

\begin{center}
\begin{tabular}{@{}lcccccc@{}}
\toprule
base $a$ & $2$ & $2.5$ & $2.7$ & $2.718$ & $2.8$ & $3$ \\
\midrule
starting slope $M(a)$ & $0.693$ & $0.916$ & $0.993$ & $1.0000$ & $1.030$ & $1.099$ \\
\bottomrule
\end{tabular}
\end{center}

Look at the trend. At base $2$ the curve crosses $(0,1)$ too \emph{shallowly} -- slope $0.69$, less than $1$. At base $3$ it crosses too \emph{steeply} -- slope $1.10$, more than $1$. Somewhere in between, the slope passes through exactly $1$. The table says it happens right around $a=2.718$. We have trapped the special base between $2$ and $3$, and pinned it near $2.718$, purely by asking the calculus question ``which curve leaves $(0,1)$ at $45^\circ$?''

\begin{figure}[htbp]
\centering
\begin{tikzpicture}
\begin{axis}[
  width=12cm, height=7cm,
  title style={font=\small},
  xlabel={$x$}, ylabel={$y$},
  domain=-1.3:1.3, samples=120,
  ymin=0, ymax=3.6, xmin=-1.4, xmax=1.45,
  axis lines=middle,
  ytick={2, 3},
  legend pos=south east, legend cell align=left, legend style={font=\footnotesize},
]
\addplot[curveblue, very thick] {2^x};
\addlegendentry{$y=2^x$ (too shallow at $0$)}
\addplot[structurepurple, very thick] {3^x};
\addlegendentry{$y=3^x$ (too steep at $0$)}
\addplot[targetgreen, very thick] {e^x};
\addlegendentry{$y=e^x$ (slope exactly $1$ at $0$)}
\addplot[tangentred, thick, dashed, domain=-1.1:1.1] {1+x};
\addlegendentry{slope-$1$ line $y=1+x$}
\node[circle,fill=black,inner sep=1.2pt] at (axis cs:0,1){};
\node[anchor=south east, font=\footnotesize] at (axis cs:-0.03,1.05) {$(0,1)$};
\end{axis}
\end{tikzpicture}
\caption{The hunt for the natural base. As $x$ leaves $0$, the curve $2^x$ rises too slowly and $3^x$ too fast. The base $e$ is the one whose tangent at $(0,1)$ is exactly the line $y=1+x$ with slope $1$.}
\label{ch03:fig-slope-one}
\end{figure}

This is our first sighting of $e$, and notice how we got it: not by decree, but by \emph{demanding} a clean derivative. The number had to exist, because the starting slope $M(a)$ slides continuously from below $1$ (at $a=2$) to above $1$ (at $a=3$), so by the Intermediate Value Theorem it equals $1$ somewhere in between. That ``somewhere'' is $e$.

\begin{definition}[The number $e$, slope view]\label{ch03:def-e-slope}
$e$ is the unique positive base for which the exponential $e^x$ has slope $1$ at $x=0$; equivalently, the unique base making
\[
\frac{\dd}{\dd x}\,e^{x} = e^{x}.
\]
Up to a constant multiple, the functions that equal their own derivative are precisely $C e^x$ (with $e^x$ being the unique one with $f(0)=1$).
\end{definition}

\begin{aestheticbox}
Pause on how special Definition~\ref{ch03:def-e-slope} is. ``A function that is its own derivative'' sounds almost paradoxical -- the curve's height \emph{is} its steepness, everywhere, simultaneously. Picture it: wherever $e^x$ is at height $5$, it is climbing at rate $5$; at height $100$, climbing at rate $100$. The function carries its own velocity inside its position. This is why $e^x$ is the natural language of every process whose speed is set by its present state, and why differential equations -- the mathematics of change -- are written in base $e$ and almost never in base $10$.
\end{aestheticbox}

\section{Three Faces of e}

We now have one definition of $e$ (the slope-$1$ base). But $e$ is famous precisely because it refuses to be just one thing. The same number arises from money, from slopes, and from a tidy infinite sum, and the fact that these all agree is the real evidence that $e$ is fundamental. We will meet \emph{three faces} of $e$ and show they are the same number.

\subsection{First face: the compound-interest limit}

Forget calculus for a moment and think about money. You invest \$1 at $100\%$ annual interest. If interest is paid once, at year's end, you finish with $\$1\times(1+1)=\$2$. Fair enough.

But suppose the bank compounds \emph{twice} a year: it pays $50\%$ at the half-year, and then $50\%$ on the new, larger balance. After six months you have $1\times(1+\tfrac12)=1.5$; after the full year, $1.5\times(1+\tfrac12)=(1+\tfrac12)^2 = 2.25$. By splitting the year you earned interest \emph{on your interest}, and beat \$2.

Why not compound monthly? Quarterly? Daily? Every second? If we split the year into $n$ equal periods, each period pays interest $1/n$, and the year-end balance is
\[
\Big(1+\frac1n\Big)^{n}.
\]
Let us tabulate this as $n$ grows.

\begin{center}
\begin{tabular}{@{}lll@{}}
\toprule
compounding & $n$ & $\left(1+\tfrac1n\right)^n$ \\
\midrule
yearly      & $1$        & $2.00000$ \\
twice       & $2$        & $2.25000$ \\
quarterly   & $4$        & $2.44141$ \\
monthly     & $12$       & $2.61304$ \\
daily       & $365$      & $2.71457$ \\
hourly      & $8760$     & $2.71813$ \\
each minute & $525600$   & $2.71828$ \\
$\to$ instant & $\infty$ & $2.718281828\ldots$ \\
\bottomrule
\end{tabular}
\end{center}

\begin{intuitionbox}
Look at the tug-of-war inside $(1+1/n)^n$ as $n$ grows: the base $1+1/n$ shrinks toward $1$ (which pulls the result toward $1$), while the exponent $n$ shoots toward $\infty$ (which pushes the result upward). Neither side wins outright; they balance at a finite ceiling near $2.718$. Going from yearly to monthly compounding makes a real jump, but going from hourly to every second barely moves the fifth decimal place.
\end{intuitionbox}

\begin{definition}[The number $e$, compound-interest view]\label{ch03:def-e-limit}
\[
e = \lim_{n\to\infty}\Big(1+\frac1n\Big)^{n}.
\]
\end{definition}

This limit really does converge -- the values are increasing but bounded above (you can check $(1+1/n)^n < 3$ for all $n$), so by the completeness of the reals they close in on a definite limit. That limit is $e \approx 2.71828$.

\begin{figure}[htbp]
\centering
\begin{tikzpicture}
\begin{axis}[
  width=12cm, height=6.4cm,
  title style={font=\small},
  xlabel={number of compounding periods $n$}, ylabel={$(1+1/n)^n$},
  xmode=log, log basis x=10,
  xmin=1, xmax=1e6, ymin=1.9, ymax=2.85,
  axis lines=left, grid=both, grid style={gray!18},
  legend pos=south east,
  legend cell align=left, legend style={font=\footnotesize},
]
\addplot[only marks, curveblue, mark=*, mark size=1.6pt] coordinates {
 (1,2.0)(2,2.25)(4,2.4414)(12,2.6130)(365,2.7146)(8760,2.7181)(525600,2.71828)};
\addlegendentry{$(1+1/n)^n$}
\addplot[tangentred, thick, dashed, domain=1:1e6] {2.718281828};
\addlegendentry{$y=e\approx2.71828$}
\end{axis}
\end{tikzpicture}
\caption{Compounding more and more finely. The balance climbs toward a hard ceiling: the horizontal line $y=e$. The horizontal axis is logarithmic, so even multiplying the compounding frequency by a thousand barely moves the dots once $n$ is large.}
\label{ch03:fig-compound}
\end{figure}

\begin{applicationbox}
\textbf{Continuous growth.} Replace ``$100\%$ for one year'' with a rate $r$ over time $t$, compounded $n$ times. The balance is $\big(1+\tfrac{rt}{n}\big)^{n}$, and as $n\to\infty$ this becomes $e^{rt}$ (substitute $m=n/(rt)$ to recover Definition~\ref{ch03:def-e-limit}). This is the \emph{continuous compounding} formula $A(t)=A_0 e^{rt}$, and it is the same formula that governs bacterial growth, radioactive decay (with $r<0$), and Newton's law of cooling. The bank's idealized ``interest paid every instant'' is mathematically identical to a population reproducing continuously. $e$ is what self-reinforcing growth looks like when you stop chopping time into pieces and let it flow.
\end{applicationbox}

\subsection{Second face: the base whose derivative is itself}

This is Definition~\ref{ch03:def-e-slope}, the slope-$1$ base. We have already cornered it numerically near $2.718$ in Figure~\ref{ch03:fig-slope-one}. The thrilling part is that it gives the \emph{same} $2.718\ldots$ as the money limit. Let us see why these two definitions give the exact same number.

\begin{whybox}
\textbf{Why the slope base equals the compound-interest limit.} Suppose $a$ is the base with starting slope $1$, so for tiny $h$,
\[
\frac{a^{h}-1}{h}\approx 1, \quad\text{i.e.}\quad a^{h}\approx 1+h .
\]
Now build up $a^1=a$ by taking $h=1/n$ and stacking $n$ of these little steps. The slope-$1$ approximation $a^{h}\approx 1+h$ with $h=1/n$ reads $a^{1/n}\approx 1+\tfrac1n$, and this gets better the larger $n$ is (the smaller the step). Since $a = (a^{1/n})^n$,
\[
a \;=\; \big(a^{1/n}\big)^{n} \;\approx\; \Big(1+\frac1n\Big)^{n}.
\]
As $n\to\infty$ the approximation tightens to equality, and the right-hand side is exactly Definition~\ref{ch03:def-e-limit}. So the base with slope $1$ \emph{is} $\lim_{n\to\infty}(1+1/n)^n$. The ``money number'' and the ``slope number'' were never two numbers. The fine-grained compounding $(1+1/n)^n$ is literally the process of building self-proportional growth out of tiny equal multiplicative steps -- which is the very thing a derivative of $1$ at the origin demands.
\end{whybox}

\begin{labbox}
Confirm the bridge with your own hands. Take the slope-$1$ condition $a^h\approx 1+h$ at $h=0.001$, which says the special base should satisfy $a \approx (1+0.001)^{1000}$. Compute $(1.001)^{1000}$ on a calculator. You should get about $2.7169$ -- already closing in on $2.71828$. Now do it again with $h=0.0001$, i.e. $(1.0001)^{10000}\approx 2.7181$. Watch the digits of $e$ assemble themselves as $h$ shrinks.
\end{labbox}

\subsection{Third face: the infinite sum $\sum 1/n!$}

The third face is the one a computer actually uses. It previews the theory of \textbf{infinite series} (which we will explore systematically in Chapter 29). For now, think of an infinite sum simply as a sequence of running totals, called \textbf{partial sums}, that draw closer and closer to a final destination:
\[
e \;=\; 1 + 1 + \frac{1}{2!} + \frac{1}{3!} + \frac{1}{4!} + \cdots
   \;=\; \sum_{n=0}^{\infty} \frac{1}{n!},
\]
where $n! = 1\cdot 2\cdot 3\cdots n$ is the \textbf{factorial}---the product of all positive integers from $1$ up to $n$, with $0!=1$ by convention. The terms decrease rapidly, because each successive term divides by an increasing integer:
\[
1,\ 1,\ \tfrac{1}{2}=0.5,\ \tfrac{1}{6}\approx0.1667,\ \tfrac{1}{24}\approx0.0417,\ \tfrac{1}{120}\approx0.0083,\ \tfrac{1}{720}\approx0.0014,\ \ldots
\]
Add them up cumulatively:
\begin{center}
\begin{tabular}{@{}lcccccc@{}}
\toprule
terms used & $2$ & $3$ & $4$ & $5$ & $6$ & $7$ \\
\midrule
partial sum & $2.0000$ & $2.5000$ & $2.6667$ & $2.7083$ & $2.7167$ & $2.7181$ \\
\bottomrule
\end{tabular}
\end{center}
After only seven terms we already have $2.7181$, correct to four decimal places. Compare that with the compound-interest limit, which needed over half a million compounding periods to reach the same precision. \emph{This} is why the series is the practical face: it converges quickly.

\begin{figure}[htbp]
\centering
\begin{tikzpicture}
\begin{axis}[
  width=12cm, height=6.2cm,
  title style={font=\small},
  xlabel={number of terms}, ylabel={partial sum},
  xmin=0.5, xmax=8.5, ymin=1.8, ymax=2.85,
  axis lines=left, grid=both, grid style={gray!18},
  legend style={at={(0.96,0.12)}, anchor=south east, font=\footnotesize},
  legend cell align=left,
  xtick={1,2,3,4,5,6,7,8},
]
\addplot[only marks, targetgreen, mark=*, mark size=1.7pt] coordinates {
 (1,1.0)(2,2.0)(3,2.5)(4,2.66667)(5,2.70833)(6,2.71667)(7,2.71806)(8,2.71825)};
\addlegendentry{partial sums of $\sum 1/n!$}
\addplot[tangentred, thick, dashed, domain=0.5:8.5] {2.718281828};
\addlegendentry{~$y=e$}
\end{axis}
\end{tikzpicture}
\caption{The series $\sum 1/n!$ races to $e$. Each new term is far smaller than the last, so the partial sums lock onto the limit after only a handful of steps -- a stark contrast to the crawling convergence of $(1+1/n)^n$.}
\label{ch03:fig-series}
\end{figure}

Where does the series come from? It is hiding inside the compound-interest limit. If you expand $(1+1/n)^n$ with the binomial theorem and let $n\to\infty$, each binomial term turns into a $1/k!$. Here is the mechanism.

\begin{whybox}
\textbf{Why $(1+1/n)^n$ unfolds into $\sum 1/k!$.} The binomial theorem gives
\[
\Big(1+\frac1n\Big)^{n} = \sum_{k=0}^{n}\binom{n}{k}\frac{1}{n^{k}}
= \sum_{k=0}^{n}\frac{n(n-1)(n-2)\cdots(n-k+1)}{k!}\cdot\frac{1}{n^{k}}.
\]
Look at the $k$-th term. The numerator $n(n-1)\cdots(n-k+1)$ is a product of $k$ factors, each just a hair below $n$; dividing by $n^{k}$ pairs them up as
\[
\frac{n}{n}\cdot\frac{n-1}{n}\cdot\frac{n-2}{n}\cdots\frac{n-k+1}{n}
= 1\cdot\Big(1-\tfrac1n\Big)\Big(1-\tfrac2n\Big)\cdots\Big(1-\tfrac{k-1}{n}\Big).
\]
As $n\to\infty$ every one of those factors heads to $1$, so the $k$-th term collapses to exactly $\dfrac{1}{k!}$. Summing over $k$,
\[
e=\lim_{n\to\infty}\Big(1+\frac1n\Big)^n = \sum_{k=0}^{\infty}\frac{1}{k!}.
\]
Rather than a separate definition, the series is what the compound-interest limit \emph{becomes} when expanded term by term; all three expressions define the same number.
\end{whybox}

The series view will return with full force in Chapter 30, when we discover that $e^x$ itself has the Taylor series $\sum x^n/n!$ -- and that this is exactly the function that equals its own derivative (differentiate term by term: $\ddx{x^n/n!} = x^{n-1}/(n-1)!$, and the sum reproduces itself, shifted). All three faces are then revealed as the same creature seen from three sides.

\begin{historybox}
The symbol $e$ and the number's role in analysis are due to Leonhard Euler in the 1720s--1740s.
The constant had surfaced earlier when Jacob Bernoulli studied continuous compound interest in 1683,
showing that $\lim_{n\to\infty}(1+1/n)^n$ lies between $2$ and $3$. Euler named the constant $e$,
computed it to $23$ decimal places using the series $\sum 1/n!$, proved it is irrational, and linked it
to trigonometry through $e^{i\theta} = \cos\theta + i\sin\theta$. In 1873, Charles Hermite proved that
$e$ is transcendental (it is never the root of a polynomial with integer coefficients).
\end{historybox}

\section{Why e Is Inevitable, and the Natural Log as Area}

\subsection{The natural logarithm}

The logarithm to base $e$ is so central that it gets its own name and symbol: the \textbf{natural logarithm},
\[
\ln x \;=\; \log_e x.
\]
It is the inverse of $e^x$: $\ln(e^x)=x$ and $e^{\ln x}=x$ for $x>0$. Everything we proved about logarithms holds, of course -- $\ln(uv)=\ln u+\ln v$, and so on -- but the natural log carries an extra gift inherited from its base. Because $e^x$ is its own derivative, its inverse $\ln x$ has the cleanest possible derivative of any logarithm. We will prove later that
\[
\frac{\dd}{\dd x}\,\ln x = \frac{1}{x},
\]
with no messy constant out front. (For any other base, $\frac{\dd}{\dd x}\log_a x = \frac{1}{x\ln a}$ -- the natural log is the one that washes the constant away, precisely because its base is the slope-$1$ base.)

\subsection{A preview: the natural log is an area}

That clean derivative $1/x$ provides another geometric path to $e$, one that needs no exponents at all. Here is the idea, as an intuitive preview of the integral chapters to come (specifically Chapter 20).

Consider the curve $y=1/t$ for $t>0$ -- a simple hyperbola sloping down toward the horizontal axis. Define $L(x)$ to be the accumulated \emph{geometric area} under this curve between $t=1$ and $t=x$. In standard calculus notation, this accumulation is written with an elongated `S' called an \textbf{integral}:
\[
L(x) \;=\; \int_1^x \frac{1}{t}\,\dd t \;=\; \text{(area beneath } y=1/t \text{ from } t=1 \text{ to } t=x).
\]
For now, you do not need any formal integration rules; simply picture shading the region on graph paper and counting squares. Yet this innocent area function turns out to \emph{be} the natural logarithm: $L(x)=\ln x$. And it carries the logarithm law right in its geometry: stretching the region horizontally by a factor reproduces $\ln(ab)=\ln a + \ln b$ as a pure statement about areas.

\begin{figure}[htbp]
\centering
\begin{tikzpicture}
\begin{axis}[
  width=12cm, height=6.4cm,
  title style={font=\small},
  xlabel={$t$}, ylabel={$y=1/t$},
  domain=0.5:5, samples=120,
  ymin=0, ymax=2.1, xmin=0, xmax=5,
  axis lines=left, grid=both, grid style={gray!18},
]
\addplot[draw=none, fill=areafill, opacity=0.6, domain=1:2.71828] {1/x} \closedcycle;
\addplot[curveblue, very thick, domain=0.45:5] {1/x};
\node at (axis cs:1.75,0.45) {\small area $=1$};
\draw[dashed, gray] (axis cs:2.71828,0) -- (axis cs:2.71828,0.3678);
\node[anchor=north] at (axis cs:2.71828,-0.02) {\small $e$};
\node[anchor=north] at (axis cs:1,-0.02) {\small $1$};
\end{axis}
\end{tikzpicture}
\caption{The natural logarithm as accumulated area under $y=1/t$. The shaded region from $t=1$ to $t=e$ has area exactly $1$. This is yet another way to single out $e$: it is the upper limit at which the running area under $1/t$ first reaches $1$, i.e. the $x$ with $\ln x = 1$.}
\label{ch03:fig-area-log}
\end{figure}

\begin{figure}[htbp]
\centering
\begin{tikzpicture}
  % Four Faces of e architectural lanes (art_costs visual architecture)
  \node[draw=sourceblue!70, fill=sourceblue!12, rounded corners=3pt, thick,
        text width=13.0cm, align=left, inner sep=6pt] (lane1) at (0, 6.45) {
    \textbf{\footnotesize FACE 1 (GROWTH \& FINANCE):}\quad
    $\displaystyle e = \lim_{n\to\infty}\left(1 + \frac{1}{n}\right)^{\!n} \approx 2.71828\dots$\\[0.55em]
    \footnotesize\textit{Continuous compounding limit: instantaneous proportional growth at all moments.}
  };

  \node[draw=targetgreen!70, fill=targetgreen!12, rounded corners=3pt, thick,
        text width=13.0cm, align=left, inner sep=6pt] (lane2) at (0, 4.30) {
    \textbf{\footnotesize FACE 2 (CALCULUS \& TANGENTS):}\quad
    $\displaystyle \frac{\dd}{\dd x} e^x = e^x \quad \Longleftrightarrow \quad \lim_{h\to 0}\frac{e^h - 1}{h} = 1$\\[0.55em]
    \footnotesize\textit{Self-derivative: curve leaves $(0,1)$ at $45^\circ$; slope equals current height.}
  };

  \node[draw=quotientorange!70, fill=quotientorange!12, rounded corners=3pt, thick,
        text width=13.0cm, align=left, inner sep=6pt] (lane3) at (0, 2.15) {
    \textbf{\footnotesize FACE 3 (SERIES EXPANSION):}\quad
    $\displaystyle e = \sum_{k=0}^\infty \frac{1}{k!} = 1 + 1 + \frac{1}{2} + \frac{1}{6} + \frac{1}{24} + \dots$\\[0.55em]
    \footnotesize\textit{Rapid factorial convergence: locks onto precision in only a handful of terms.}
  };

  \node[draw=structurepurple!70, fill=structurepurple!12, rounded corners=3pt, thick,
        text width=13.0cm, align=left, inner sep=6pt] (lane4) at (0, 0) {
    \textbf{\footnotesize FACE 4 (GEOMETRIC ACCUMULATION):}\quad
    $\displaystyle \int_1^e \frac{1}{t}\,\dd t = 1 \quad \Longleftrightarrow \quad \ln(e) = 1$\\[0.55em]
    \footnotesize\textit{Hyperbolic area: boundary where accumulated area under $1/t$ reaches unit area.}
  };
\end{tikzpicture}
\caption{Four distinct doors opening onto the exact same mathematical constant $e \approx 2.71828\dots$. Whether arrived at through continuous compound interest, the requirement that an exponential equal its own derivative, the binomial expansion into reciprocal factorials, or the natural accumulation under $y=1/t$, the identical universal number appears.}
\label{ch03:fig-four-faces}
\end{figure}

\begin{aestheticbox}
Why is $e$ called ``natural'' while base $10$ is not? Base $10$ is an accident of human anatomy; base $e$ is forced by calculus itself. Ask for continuous growth, ask for an exponential whose derivative needs no scaling constant, expand that growth into factorials, or measure the area under $1/t$ -- in each case, $e$ is the value you find.
\end{aestheticbox}

\begin{connectionbox}
This area-under-$1/t$ picture is more than a curiosity -- it is the cleanest \emph{logical} foundation for the whole subject. Notice the awkwardness we waded through in the opening section: to define $a^x$ for irrational $x$ we needed limits, and to differentiate it we needed the mysterious constant $M(a)$. The modern, airtight approach turns the whole story upside down. \emph{First} define $\ln x$ as the area $\int_1^x \frac{1}{t}\,\dd t$ -- which needs only integration, no irrational exponents at all. \emph{Then} define $e^x$ as its inverse, and \emph{finally} define $a^x = e^{x\ln a}$. Every property we fought for -- the exponent law, the derivative rule, the value of $e$ -- drops out as a clean theorem. We will walk that road in the chapters on the integral and the Fundamental Theorem. For now, the hyperbola in Figure~\ref{ch03:fig-area-log} is quietly holding up everything in this chapter.
\end{connectionbox}

\section{Worked Examples}

\begin{example}[Reading the exponent law backward]\label{ch03:ex-law}
Simplify $\dfrac{e^{3}\cdot e^{x}}{e^{x-2}}$.

\textbf{Solution.} Combine the numerator with the law $e^a e^b = e^{a+b}$: the top is $e^{3+x}$. Dividing by $e^{x-2}$ subtracts exponents:
\[
\frac{e^{3+x}}{e^{x-2}} = e^{(3+x)-(x-2)} = e^{3+x-x+2} = e^{5}.
\]
The $x$'s cancel entirely. Whenever you see a ratio of exponentials with the same base, do not panic -- just add and subtract the exponents.
\end{example}

\begin{example}[Logs turn an equation inside out]\label{ch03:ex-solve}
Solve $5\cdot 2^{x} = 320$ for $x$.

\textbf{Solution.} First isolate the exponential: $2^x = 64$. Now ask the logarithm's question, ``$2$ to the what is $64$?'' Since $64=2^6$, we read off $x=6$. If the right-hand side were not a clean power of $2$, we would take $\ln$ of both sides:
\[
x\ln 2 = \ln 64 \quad\Longrightarrow\quad x = \frac{\ln 64}{\ln 2} = \log_2 64 = 6 .
\]
The logarithm is exactly the tool that brings the unknown \emph{down} from the exponent where you cannot reach it.
\end{example}

\begin{example}[Continuous vs. annual compounding]\label{ch03:ex-money}
You invest \$1000 at $6\%$ per year for $10$ years. Compare yearly compounding with continuous compounding.

\textbf{Solution.} Yearly: $A = 1000(1+0.06)^{10} = 1000(1.06)^{10}\approx 1000(1.7908)=\$1790.85$. Continuous: $A = 1000\,e^{0.06\cdot 10} = 1000\,e^{0.6}\approx 1000(1.8221)=\$1822.12$. Compounding ``every instant'' instead of once a year earns an extra \$31 over the decade -- real money, but a far cry from doubling. The gap is the diminishing return we saw in Figure~\ref{ch03:fig-compound}: most of the benefit of compounding is captured already by compounding a few times a year.
\end{example}

\begin{example}[How many terms of the series?]\label{ch03:ex-series}
Estimate $e$ using the series through the $1/4!$ term, and bound the error.

\textbf{Solution.} The partial sum is
\[
1 + 1 + \tfrac12 + \tfrac16 + \tfrac{1}{24} = 2 + 0.5 + 0.16667 + 0.04167 = 2.70833.
\]
The true value is $2.71828$, so the error is about $0.00995$. The first omitted term is $1/5! = 1/120 \approx 0.00833$, and indeed the total error is just a touch more than that one term -- because every \emph{later} term is far smaller still. This is the hallmark of a fast series: the error is essentially the size of the first thing you dropped.
\end{example}

\begin{example}[Estimating the starting slope of $3^x$]\label{ch03:ex-slope}
Estimate $M(3)=\lim_{h\to0}(3^h-1)/h$ numerically, and say what it means.

\textbf{Solution.} With $h=0.001$: $3^{0.001}=e^{0.001\ln 3}\approx 1+0.001(1.0986)=1.0010986$, so $(3^{0.001}-1)/0.001 \approx 1.0986$. Thus $M(3)\approx 1.099$, matching our table. It means the curve $3^x$ crosses $(0,1)$ at slope $1.099$ -- a bit steeper than $45^\circ$ -- which is why $3$ overshoots the natural base $e$. (And note $1.0986\approx \ln 3$: the starting slope of $a^x$ is precisely $\ln a$, a fact we will prove once we have the chain rule, since $a^x = e^{x\ln a}$.)
\end{example}

\section{Looking Ahead: Periodic Change}
\label{ch03:looking-ahead}

Exponentials model growth and decay, where quantities multiply and move away from their starting values.
Physical quantities also cycle and oscillate: pendulums swing, strings vibrate, and planets orbit.
In Chapter~4, we turn to the functions that describe periodic change: sines and cosines.

\section{Exercises}

\begin{exercise}[Proof: algebraic laws of exponents]
\label{ch03:exc-laws}
Using only the exponent law $a^{x+y}=a^x a^y$ and $a^0=1$, prove that $a^{-x}=1/a^x$ and that $(a^x)^y = a^{xy}$ for integer $y$. State clearly where you use each law.
\end{exercise}

\begin{exercise}[Mental check: evaluating logarithms]
\label{ch03:exc-eval-logs}
Simplify without a calculator: (a) $\log_2 32$; (b) $\log_3 \tfrac{1}{27}$; (c) $\ln(e^4)$; (d) $e^{\ln 7}$; (e) $\log_5 1$. For each, say in one sentence what question the logarithm is answering.
\end{exercise}

\begin{exercise}[Solving exponential and logarithmic equations]
\label{ch03:exc-solve-eqs}
Solve for $x$: (a) $3^{x}=81$; (b) $4\cdot 10^{x}=4000$; (c) $e^{2x}=5$ (leave the answer in terms of $\ln$); (d) $\ln(x-1)=2$.
\end{exercise}

\begin{exercise}[Application: bacterial growth and continuous rate]
\label{ch03:exc-bacteria}
A culture of bacteria doubles every $3$ hours and starts at $500$ cells. (a) Write the population as $N(t)=500\cdot 2^{t/3}$ and explain why the exponent is $t/3$. (b) How many cells after $12$ hours? (c) Rewrite the model in the form $N(t)=500\,e^{kt}$ by finding $k$, and interpret $k$ as a continuous growth rate.
\end{exercise}

\begin{exercise}[Mental check: doubling time and half-life]
\label{ch03:exc-doubling-half}
A radioactive sample decays according to $N(t) = N_0 2^{-t/T_{\text{half}}}$, where $T_{\text{half}}$ is its half-life.
\begin{enumerate}[label=(\alph*)]
  \item Explain in words why every tick of $T_{\text{half}}$ units of time halves the remaining mass.
  \item What fraction of the original sample remains after $3$ half-lives? After $5$ half-lives?
  \item Convert this decay formula into continuous base-$e$ form, $N(t) = N_0 e^{-\lambda t}$, and find the decay constant $\lambda$ in terms of $T_{\text{half}}$ and $\ln 2$.
\end{enumerate}
\end{exercise}

\begin{exercise}[Squeezing e: compound interest bounds]
\label{ch03:exc-compound-bounds}
Compute $(1+1/n)^n$ by hand for $n=1,2,3,4$ and confirm the sequence is increasing. Then compute $(1+1/n)^{n+1}$ for the same $n$ and confirm \emph{that} sequence is decreasing. (These two sequences squeeze $e$ from below and above -- a clean way to see the limit exists.)
\end{exercise}

\begin{exercise}[Series computation and error bound]
\label{ch03:exc-series-error}
Using the series $e=\sum_{n=0}^{\infty}1/n!$, compute the partial sum through $1/6!$ and compare with $2.718282$. How many decimal places do you get right, and how does the error compare to the first omitted term $1/7!$?
\end{exercise}

\begin{exercise}[Conceptual: slope equals height]
\label{ch03:exc-conc-slope}
Explain in a short paragraph, without formulas, why a function whose derivative equals itself must grow exponentially. (Hint: at every height its steepness equals its height, so the bigger it gets, the faster it climbs -- describe the runaway feedback.)
\end{exercise}

\begin{exercise}[Conceptual: what makes e natural?]
\label{ch03:exc-conc-natural}
Why is base $e$ called ``natural'' while base $10$ is merely convenient for humans? Give two distinct reasons drawn from this chapter -- one about derivatives and one about the area under $1/t$.
\end{exercise}

\begin{exercise}[Change of base as vertical stretching]
\label{ch03:exc-change-base}
Show, using the change-of-base formula, that the graphs of $\log_2 x$, $\ln x$, and $\log_{10}x$ are all vertical rescalings of one another. By what factor must you stretch $\ln x$ to obtain $\log_{10}x$?
\end{exercise}

\begin{exercise}[Proof: properties of the starting slope]
\label{ch03:exc-proof-slope}
Prove that the starting slope $M(a)=\lim_{h\to0}(a^h-1)/h$ satisfies $M(a^2)=2M(a)$, and more generally $M(a^k)=k\,M(a)$ for positive integers $k$. (Hint: $a^{2h}-1 = (a^h-1)(a^h+1)$, and $a^h\to 1$ as $h\to 0$.) Explain why this is consistent with the claim that $M(a)=\ln a$.
\end{exercise}

\begin{exercise}[Connection: hyperbolic functions and the circle preview]
\label{ch03:exc-hyperbolic}
Recall from Chapter~2 (Exercise~\ref{ch02:exc-proof-decomp}) that any function whose domain is symmetric about $0$ can be uniquely split into an even part and an odd part. For the natural exponential $f(x) = e^x$, this decomposition yields the \textbf{hyperbolic cosine} and \textbf{hyperbolic sine}:
\[
\cosh x \;=\; \frac{e^x + e^{-x}}{2}, \qquad \sinh x \;=\; \frac{e^x - e^{-x}}{2}.
\]
\begin{enumerate}[label=(\alph*)]
  \item Expand the squares algebraically and prove the fundamental hyperbolic identity:
  \[
  (\cosh x)^2 - (\sinh x)^2 \;=\; 1.
  \]
  \item What geometric curve in the $(u,v)$-plane does the parameterized point $(u,v) = (\cosh x, \sinh x)$ trace out as $x$ varies? (Compare the hyperbolic curve $u^2 - v^2 = 1$ with the unit circle $u^2 + v^2 = 1$ that anchors Chapter~4.)
  \item Confirm that $\cosh 0 = 1$ and $\sinh 0 = 0$.
\end{enumerate}
\end{exercise}

\begin{exercise}[Challenge: Fourier's irrationality proof for e]
\label{ch03:exc-challenge-fourier}
The number $e$ is irrational. Here is Fourier's argument. Suppose for contradiction $e=p/q$ for integers $p,q$. Multiply the series $e=\sum 1/n!$ by $q!$. Argue that $q!\,e$ would have to be an integer, but that the tail $\sum_{n=q+1}^{\infty} q!/n!$ is a positive number strictly less than $1$ (compare it with a geometric series). Conclude that no such $p,q$ exist. Fill in the comparison that bounds the tail below $1$.
\end{exercise}
